\documentclass[12pt,a4paper]{article}

% Documento académico monocromático, tipografía Times New Roman.
% Si Times New Roman no está instalada, Tinos conserva una apariencia
% y métricas muy cercanas. XeLaTeX selecciona automáticamente la fuente.
\usepackage[spanish,es-noshorthands]{babel}
\usepackage{fontspec}
\IfFontExistsTF{Times New Roman}{%
  \setmainfont{Times New Roman}%
}{%
  \setmainfont{Tinos}%
}
\usepackage[a4paper,top=2.5cm,bottom=2.5cm,left=2.6cm,right=2.6cm]{geometry}
\usepackage{setspace}
\setstretch{1.18}
\usepackage{array}
\usepackage{tabularx}
\usepackage{booktabs}
\usepackage{ragged2e}
\usepackage{titlesec}
\usepackage{fancyhdr}
\usepackage{needspace}
\usepackage{microtype}
\usepackage{etoolbox}

% Presentación completamente negra.
\usepackage{xcolor}
\color{black}
\AtBeginDocument{\color{black}}

% Jerarquía académica sin colores ni elementos decorativos.
\titleformat{\section}{\normalfont\bfseries\fontsize{12}{15}\selectfont}{\thesection.}{0.5em}{}
\titleformat{\subsection}{\normalfont\bfseries\fontsize{12}{15}\selectfont}{}{0pt}{}
\titlespacing*{\section}{0pt}{15pt}{7pt}
\titlespacing*{\subsection}{0pt}{11pt}{4pt}
\setlength{\parindent}{0pt}
\setlength{\parskip}{5pt}
\renewcommand{\arraystretch}{1.25}
\setlength{\tabcolsep}{5pt}
\pagestyle{fancy}
\fancyhf{}
\renewcommand{\headrulewidth}{0pt}
\fancyfoot[C]{\footnotesize\thepage}

\begin{document}
\begin{center}
  {\bfseries\fontsize{14}{17}\selectfont ACTA DE ENTREVISTA}\par
  \vspace{3pt}
  {\fontsize{12}{15}\selectfont Levantamiento de requisitos -- Sistema de Gestión de Ayudantías de Cátedra}
\end{center}
\vspace{6pt}

\noindent\begin{tabularx}{\textwidth}{@{}>{\bfseries}p{4.8cm}X@{}}
Proyecto: & TA-PFC-E1 -- Sistema de Gestión de Ayudantías de Cátedra (SGAC) \\
Fecha de entrevista: & 21-11-2026 \\
Fecha de revisión documental: & 22-11-2026 \\
Medio: & Discord \\
Entrevistado: & KHRIZ ALEXANDER CORONEL GOMEZ \\
Cargo: & Ayudante de Cátedra \\
Participantes: & MUÑOZ VERA MELANIE MICHELLE y SANCHEZ VERA DARWIN JOEL \\
Responsable de documentación: & QUINTERO RICO LUIS DAVID \\
\end{tabularx}

\section{Objetivo}
Recopilar información sobre el proceso de postulación y gestión de ayudantías de cátedra, sus dificultades y las necesidades que deben considerarse para el SGAC.

\section{Síntesis de temas y necesidades}
Los identificadores se relacionan con el guion de entrevista. Se presenta una síntesis temática para el análisis de requisitos, no una transcripción literal.

\subsection{P10. Volumen de postulaciones}
Para dimensionar la plataforma se requiere conocer el número de postulaciones por convocatoria. El equipo considera \textbf{350 postulaciones} como escenario inicial de diseño, aún sin respaldo estadístico institucional.

\subsection{P12. Usuarios simultáneos}
La demanda podría concentrarse cerca del cierre de postulaciones y de la publicación de resultados. Se plantean escenarios de prueba de \textbf{50 usuarios simultáneos} en consultas y \textbf{100 usuarios} en momentos de mayor actividad.

\subsection{P13. Tiempos de respuesta}
La consulta del estado de postulación debe ser ágil. Se plantea como objetivo técnico inicial un \textbf{tiempo de respuesta p95 inferior a 300 ms}, bajo una carga de 50 usuarios simultáneos.

\subsection{P14. Disponibilidad}
La disponibilidad es especialmente importante durante el período de inscripción. Se propone inicialmente un objetivo del \textbf{99\,\% durante cinco días de recepción de postulaciones}.

\subsection{P15. Consulta de información académica}
Se contempla consultar en el SGA la matrícula, las calificaciones y el promedio del postulante para verificar su elegibilidad, mediante una integración de solo lectura. Su viabilidad requiere revisión técnica.

\subsection{P16. Fallos del SGA}
Cuando el SGA no esté disponible, la comprobación automática de los requisitos académicos no podrá concluirse de inmediato. El procedimiento institucional de contingencia aún debe precisarse.

\subsection{P17. Manejo de postulaciones durante fallos del SGA}
Se propone conservar las solicitudes válidas con el estado \textbf{«Validación pendiente»} y reintentar la consulta al restablecerse el SGA, evitando duplicaciones o nuevas postulaciones innecesarias.

\subsection{P21. Respaldos y recuperación}
Para proteger los expedientes y las evidencias de ayudantía se proponen inicialmente \textbf{respaldos cada 24 horas} y \textbf{recuperación en un máximo de cuatro horas}, sujetos a evaluación técnica.

\section{Relación con los requisitos no funcionales}
\noindent\begin{tabularx}{\textwidth}{@{}p{1.5cm}p{2.6cm}X@{}}
\toprule
\textbf{RNF} & \textbf{Referencia} & \textbf{Criterio técnico inicial} \\
\midrule
RNF-01 & P12, P13 & 50 usuarios simultáneos; p95 inferior a 300 ms. \\
RNF-02 & P10, P12 & 350 postulaciones por convocatoria; pico de 100 usuarios. \\
RNF-03 & P14 & 99\,\% de disponibilidad durante cinco días de inscripción. \\
RNF-04 & P16, P17 & Conservar postulaciones en «Validación pendiente» y reintentar la verificación. \\
RNF-05 & P21 & Respaldo cada 24 horas y recuperación en un máximo de cuatro horas. \\
\bottomrule
\end{tabularx}

\section{Observaciones}
Los valores de capacidad, rendimiento, disponibilidad y recuperación son \textbf{estimaciones y metas propuestas por el equipo}, no mediciones institucionales ni declaraciones textuales atribuidas al entrevistado. Su adopción como requisitos definitivos requiere contraste con las respuestas de la entrevista y validación de los responsables académicos y técnicos.

\textbf{Revisión documental:} 22-11-2026.

\needspace{4.5cm}
\section{Firma del entrevistado}
\vspace{1.25cm}
\noindent\rule{8.5cm}{0.35pt}\par
\textbf{KHRIZ ALEXANDER CORONEL GOMEZ}\par
Firma del entrevistado

\end{document}
